\chapter{Overfitting}\label{ch:rs:overfitting}

A thousand parameter settings of a moving-average trend system are backtested on eighteen years of daily returns. The best has a Sharpe ratio
of 0.44; over the next five years it loses, with a Sharpe ratio of $-0.32$. The returns were pure noise, and the thousand systems a search through
it. Run the same search on returns with a real, planted trend and the best system in the search years again disappoints afterwards, earning less
than the average system: the trend was there for all of them, and the search picked noise on top of it. Overfitting in finance is rarely a model
with too many parameters; it is a choice among too many \omterm{def:rs:vectorised-backtests:backtest}{backtests}, measured on too little history. This chapter measures it (the \omterm{def:rs:overfitting:pbo}{probability of
backtest overfitting}), bounds it (the \omterm{def:rs:overfitting:minbtl}{minimum backtest length}), and prevents the leak that makes it invisible in cross-validation (\omterm{def:rs:overfitting:purging}{purging} and
\omterm{def:rs:overfitting:purging}{embargo}), with \texttt{firm.overfit}.

\section{Selection under many trials}

\begin{definition}[Backtest overfitting, in-sample, out-of-sample]\label{def:rs:overfitting:overfitting}
\emph{In-sample}\index{in-sample} data are those used to build or choose a strategy; \emph{out-of-sample}\index{out-of-sample} data are those
it has not seen. \emph{Backtest overfitting}\index{backtest overfitting} is the selection of a strategy whose in-sample performance owes more to
the particular sample than to a property that persists, so that its out-of-sample performance falls short of what was selected.
\end{definition}

The chapter's systems hold the market long when a fast moving average of the price is above a slow one, short otherwise: fast windows of 2 to 40
days, slow windows of 50 to 540, a thousand pairs. On 25 years of synthetic daily returns, the first 20 are the search (17.9 years once the longest
average has started), the last five the test. On pure noise the best system's \omterm{def:rs:overfitting:overfitting}{in-sample} Sharpe ratio is 0.44, its \omterm{def:rs:overfitting:overfitting}{out-of-sample} $-0.32$; the median
system's \omterm{def:rs:overfitting:overfitting}{in-sample} 0.17. The expected best of a thousand independent null strategies over 17.9 years would be 0.77 (Book~4, chapter~12); the systems
are far from independent (their Sharpe ratios have a standard deviation of only 0.076 around a common level), and the maximum expected from that
dispersion is 0.25 above the common level, about what was found.

With a planted trend, a persistent drift of the kind trend-followers earn (a half-life of 120 days), every system earns it: the median \omterm{def:rs:overfitting:overfitting}{in-sample} Sharpe
ratio is 0.50, the average \omterm{def:rs:overfitting:overfitting}{out-of-sample} 0.26. The best \omterm{def:rs:overfitting:overfitting}{in-sample} system, 0.77, earns 0.14 out of sample, less than the average system. The ranking is not useless (across the systems, \omterm{def:rs:overfitting:overfitting}{in-sample} and
\omterm{def:rs:overfitting:overfitting}{out-of-sample} Sharpe ratios correlate at 0.33, and the fifty best in sample earn 0.29 out of sample), but the single winner is where the noise concentrates: the
search found the trend, and choosing one point of it paid for noise. On noise the correlation is $-0.05$ and the fifty best earn $-0.02$.

\section{The probability of backtest overfitting}

\begin{definition}[Probability of backtest overfitting, combinatorially symmetric cross-validation]\label{def:rs:overfitting:pbo}
The \emph{probability of backtest overfitting}\index{probability of backtest overfitting} (PBO) of a selection procedure is the probability that the
strategy it selects as the best in sample performs below the median of the candidates out of sample. \emph{Combinatorially symmetric
cross-validation}\index{combinatorially symmetric cross-validation} (CSCV) estimates it by cutting the history into $S$ blocks and, for each choice of
$S/2$ blocks as the \omterm{def:rs:overfitting:overfitting}{in-sample} set, ranking the \omterm{def:rs:overfitting:overfitting}{in-sample} winner among all candidates on the other half; the PBO is the share of splits where its rank
is below the median (Bailey, Borwein, López de Prado and Zhu).
\end{definition}

The rank $\omega$ of the winner among $N$ candidates, as a fraction, is turned into a logit $\lambda = \ln(\omega/(1 - \omega))$: positive when the winner
stays above the median, negative when it falls below. With 16 blocks there are 12\,870 splits; \texttt{firm.overfit.cscv} uses 3\,000 of them, drawn at
random, and computes every split's Sharpe ratios as matrix products of block sums. The PBO of the thousand-system search is 0.82 on noise and 0.74 with the
planted trend (\cref{fig:rs:overfitting:logits}); the regression of the winners' \omterm{def:rs:overfitting:overfitting}{out-of-sample} Sharpe ratios on their \omterm{def:rs:overfitting:overfitting}{in-sample} ones has slopes of $-0.89$
and $-0.91$. A PBO above one half means selection is worse than choosing at random: among correlated systems, the \omterm{def:rs:overfitting:overfitting}{in-sample} winner is the one that fitted
the noise of its half best, and the noise of the other half is its mirror.

\begin{omfigure}
\begin{tikzpicture}
\begin{axis}[width=11cm,height=5.4cm,ybar,bar width=7pt,xlabel={logit of the in-sample winner's out-of-sample rank},ylabel={share of splits},
  tick label style={font=\small},label style={font=\small},xmin=-6.2,xmax=6.2,ymin=0,grid=major,grid style={gray!25},
  yticklabel style={/pgf/number format/fixed,/pgf/number format/precision=2},scaled y ticks=false,legend columns=2,
  legend style={font=\footnotesize,at={(0.5,-0.3)},anchor=north,draw=none,fill=none,
  /tikz/every even column/.append style={column sep=8pt}},table/col sep=comma]
\addplot[fill=gray!45,draw=gray] table[x=x,y=share]{figdata/research/20-overfitting/logits_noise.csv};
\addplot[fill=omDef!55,draw=omDef] table[x=x,y=share]{figdata/research/20-overfitting/logits_trend.csv};
\draw[omThm,thick,dashed] (axis cs:0,0) -- (axis cs:0,0.2);
\legend{noise (PBO 0.82),planted trend (PBO 0.74)}
\end{axis}
\end{tikzpicture}

\omcaption{CSCV on the thousand trend systems' search years (16 blocks, 3\,000 splits): the distribution of the logit of the \omterm{def:rs:overfitting:overfitting}{in-sample} winner's \omterm{def:rs:overfitting:overfitting}{out-of-sample}
rank. Mass to the left of zero is the \omterm{def:rs:overfitting:pbo}{probability of backtest overfitting}. Data: \texttt{rs\_overfit} on synthetic returns.}\label{fig:rs:overfitting:logits}
\end{omfigure}

\begin{omfigure}
\begin{tikzpicture}
\begin{axis}[name=a,width=6.6cm,height=5.6cm,title={\small noise},xlabel={\small in-sample Sharpe ratio},ylabel={\small out-of-sample Sharpe ratio},
  tick label style={font=\footnotesize},xmin=-0.3,xmax=0.9,ymin=-0.8,ymax=1.2,grid=major,grid style={gray!25},table/col sep=comma]
\addplot[only marks,mark=*,mark size=0.8pt,gray] table[x=is,y=oos]{figdata/research/20-overfitting/scatter_noise.csv};
\end{axis}
\begin{axis}[at={(a.east)},anchor=west,xshift=1.2cm,width=6.6cm,height=5.6cm,title={\small planted trend},xlabel={\small in-sample Sharpe ratio},
  tick label style={font=\footnotesize},xmin=-0.3,xmax=0.9,ymin=-0.8,ymax=1.2,grid=major,grid style={gray!25},table/col sep=comma]
\addplot[only marks,mark=*,mark size=0.8pt,omDef] table[x=is,y=oos]{figdata/research/20-overfitting/scatter_trend.csv};
\end{axis}
\end{tikzpicture}

\omcaption{\omterm{def:rs:overfitting:overfitting}{In-sample} (17.9 years) against \omterm{def:rs:overfitting:overfitting}{out-of-sample} (5 years) Sharpe ratios of the trend systems (every fourth of the thousand). On noise the clouds are unrelated (correlation $-0.05$); with the trend they correlate at 0.33, yet the top \omterm{def:rs:overfitting:overfitting}{in-sample} system is not near the top out of sample. Data: \texttt{rs\_overfit.systems}.}\label{fig:rs:overfitting:scatter}
\end{omfigure}

\section{Minimum backtest length}

\begin{definition}[Minimum backtest length]\label{def:rs:overfitting:minbtl}
The \emph{minimum backtest length}\index{minimum backtest length} for $N$ trials and a target Sharpe ratio is the number of years of history below which the
expected best of $N$ strategies with no skill has an \omterm{def:rs:overfitting:overfitting}{in-sample} Sharpe ratio at or above the target.
\end{definition}

\begin{proposition}[How much history a search needs]\label{prop:rs:overfitting:minbtl}
If $N$ independent strategies with no skill are each measured over $y$ years, their annual Sharpe ratio estimates have variance about $1/y$, and the expected
maximum grows like $\sqrt{2\ln N/y}$; a target Sharpe ratio $s$ is expected to be reached by chance unless $y \gtrsim 2\ln N/s^2$.
\end{proposition}

\begin{proof}
The expected maximum of $N$ standard normals grows like $\sqrt{2\ln N}$ (Book~4, chapter~12); scale by the standard error $1/\sqrt y$ and solve
$\sqrt{2\ln N/y} = s$.
\end{proof}

With the expected maximum computed exactly (\texttt{firm.multitest.expected\_max\_sr}), a thousand trials need 10.6 years before a Sharpe ratio of 1 stops
being expected from luck, 42.4 years for 0.5, and 3.3 years for 1.8 (\cref{fig:rs:overfitting:minbtl}). The bound is conservative when the trials are
correlated, as the trend systems are, and a guide when they are not; its lesson is the scaling: four times the history for half the Sharpe ratio. Harvey and
Liu (2015) turned the same logic into a haircut for any reported Sharpe ratio, given the number of tests behind it.

\begin{omfigure}
\begin{tikzpicture}
\begin{semilogxaxis}[width=10cm,height=5.6cm,xlabel={number of trials},ylabel={years of history},tick label style={font=\small},label style={font=\small},
  xmin=2,xmax=10000,ymin=0,ymax=60,grid=major,grid style={gray!25},legend pos=north west,legend style={font=\footnotesize},table/col sep=comma]
\addplot[omThm,very thick] table[x=trials,y=sr05]{figdata/research/20-overfitting/minbtl.csv};
\addplot[omDef,very thick] table[x=trials,y=sr10]{figdata/research/20-overfitting/minbtl.csv};
\addplot[black!60,thick,dashed] table[x=trials,y=sr18]{figdata/research/20-overfitting/minbtl.csv};
\legend{target Sharpe 0.5,target Sharpe 1.0,target Sharpe 1.8}
\end{semilogxaxis}
\end{tikzpicture}

\omcaption{\omterm{def:rs:overfitting:minbtl}{Minimum backtest length}: the years of history below which the expected best of $N$ independent null strategies reaches the target Sharpe ratio.
Data: \texttt{firm.overfit.min\_backtest\_length}.}\label{fig:rs:overfitting:minbtl}
\end{omfigure}

\section{Purged and embargoed cross-validation}

\begin{definition}[Label overlap, purging, embargo, combinatorial purged cross-validation]\label{def:rs:overfitting:purging}
\emph{Label overlap}\index{label overlap} occurs when the targets of different observations share periods, as $h$-day forward returns observed every day do.
\emph{Purging}\index{purging} removes from the training set every observation whose label overlaps in time the labels of the test set; an
\emph{embargo}\index{embargo} also removes the training observations that immediately follow the test set, since features are serially correlated.
\emph{Combinatorial purged cross-validation}\index{combinatorial purged cross-validation} (CPCV) uses every choice of $k$ test folds out of $n$, purged and
embargoed, recombining the test folds into $\binom{n}{k}k/n$ \omterm{def:rs:vectorised-backtests:backtest}{backtest} paths (López de Prado, 2018; as implemented in open-source libraries such as skfolio).
\end{definition}

Shuffled $k$-fold cross-validation, the default of general machine-learning libraries, puts neighbouring days into training and test sets alike. When
labels overlap, a model can predict a test day's label from the training days whose labels share its periods. The chapter's demonstration is extreme on
purpose: 20-day forward returns of pure noise, predicted by the five nearest neighbours in a feature that moves as slowly as time (a regime variable), five
folds:

\begin{center}
\small
\begin{tabular}{@{}lcccc@{}}
\toprule
cross-validation & shuffled & contiguous & purged & purged and embargoed\\
\midrule
$R^2$ of a forecast of noise & 0.90 & $-0.63$ & $-0.57$ & $-0.93$\\
\bottomrule
\end{tabular}
\end{center}

Shuffled folds report that the model explains 90\% of the variance of noise; every scheme that respects time says it is worse than a constant. Contiguous
folds leak only at their edges, \omterm{def:rs:overfitting:purging}{purging} closes the edges, and the \omterm{def:rs:overfitting:purging}{embargo} guards against features that carry information across them. With a combinatorial
scheme the same model can be judged on many paths instead of one, which turns a single \omterm{def:rs:overfitting:overfitting}{out-of-sample} Sharpe ratio into a distribution.

\section{Walk-forward analysis}

\begin{definition}[Walk-forward analysis]\label{def:rs:overfitting:walkforward}
\emph{Walk-forward analysis}\index{walk-forward analysis} repeats, through time, the cycle of choosing a strategy on a trailing window and trading it on the
next period, and judges the stitched \omterm{def:rs:overfitting:overfitting}{out-of-sample} periods.
\end{definition}

Walk-forward is the procedure the firm would actually follow, and it measures the procedure rather than a strategy. Re-choosing the best of the thousand
systems every year on the previous five years gives a stitched Sharpe ratio of $-0.04$ on noise and 0.29 with the planted trend: the trend's true Sharpe ratio,
about what the average system earns, and twice what the \omterm{def:rs:overfitting:overfitting}{in-sample} winner of the whole search earned. Its limit is the single path: one history, one sequence of
choices; CPCV and CSCV exist because one path is a small sample.

\section{Practical defences}

The defences are organisational before they are statistical. Count every trial, in the \omterm{def:rs:the-research-process:log}{research log} (chapter~1), including those abandoned; report the PBO of the
search and the deflated Sharpe ratio of its winner (Book~4, chapter~12); fix the evaluation protocol before seeing the results; keep a holdout that nobody touches
until the decision; prefer the broad region of parameters that works to its best point (a plateau, not a peak: here, the average system beat the winner); and
demand a mechanism, because a mechanism survives the change of sample that a fitted number does not. Arnott, Harvey and Markowitz (2019) set these out as a
protocol for research with machine learning; the PBO and the \omterm{def:rs:overfitting:minbtl}{minimum backtest length} put numbers on them.

\section{Tutorial: one thousand trend systems}\label{tut:rs:overfitting:tut}

\begin{tutorial}
\textbf{Goal.} Search a thousand trend systems on noise and on a planted trend, estimate the PBO by CSCV, walk forward, compute the \omterm{def:rs:overfitting:minbtl}{minimum backtest length}, and
compare four cross-validation schemes on overlapping labels. \textbf{End state:}
\cref{fig:rs:overfitting:logits,fig:rs:overfitting:scatter,fig:rs:overfitting:minbtl}; the table.
\begin{tutsteps}
\item \textbf{CSCV}: block sums, then every split's Sharpe ratios as matrix products, then the winner's \omterm{def:rs:overfitting:overfitting}{out-of-sample} rank.
\omcode{code/firm/overfit/firm_overfit.py}{33}{64}{The probability of backtest overfitting by CSCV.}\label{lst:rs:overfitting:cscv}
\item \textbf{Purged k-fold with an \omterm{def:rs:overfitting:purging}{embargo}.}
\omcode{code/firm/overfit/firm_overfit.py}{67}{80}{Purging and embargo.}\label{lst:rs:overfitting:purge}
\item \textbf{Run} \texttt{search()} for both worlds, \texttt{walk\_forward\_sr()}, \texttt{minbtl()}, \texttt{label\_leak()} and \texttt{fig\_overfit.py}.
\end{tutsteps}
\textbf{What to change next.} Choose the centre of the best plateau of parameters instead of the single best point and compare its \omterm{def:rs:overfitting:overfitting}{out-of-sample} Sharpe ratio; run
CSCV on a hundred independent strategies instead of a thousand correlated ones and see the PBO approach one half on noise.
\end{tutorial}

\section{Build: the overfitting toolkit}\label{bld:rs:overfitting:overfit}

\begin{build}
\bfield{Purpose} Every search the firm runs is reported with its probability of overfitting, its \omterm{def:rs:overfitting:minbtl}{minimum backtest length}, and cross-validation that respects
time.
\bfield{Interface} \texttt{cscv(M, n\_blocks, max\_splits, seed)}, \texttt{purged\_kfold(start, end, k, embargo)}, \texttt{cpcv(start, end, n\_folds, n\_test, embargo)},
\texttt{cpcv\_paths}, \texttt{walk\_forward(n, train, test, step, anchored)}, \texttt{min\_backtest\_length(n\_trials, sr\_target)}.
\bfield{Rules} No shuffled cross-validation on time series; labels carry their start and end; the \omterm{def:rs:the-research-process:log}{trial count} comes from the \omterm{def:rs:the-research-process:log}{research log}, not memory.
\bfield{Acceptance tests} \texttt{code/firm/overfit/tests/}: a PBO near one half and a non-positive slope on noise, near zero with one skilled strategy; \omterm{def:rs:overfitting:purging}{purging} and
\omterm{def:rs:overfitting:purging}{embargo} by hand; CPCV's fifteen splits and five paths, disjoint train and test sets; walk-forward windows, anchored and rolling; the \omterm{def:rs:overfitting:minbtl}{minimum backtest length}'s
monotonicity and its quadrupling for half the target.
\bfield{Stretch} CPCV Sharpe-ratio distributions for a model; the deflated Sharpe ratio from the effective number of trials; a PBO dashboard fed by the \omterm{def:rs:the-research-process:log}{research log}.
\end{build}

\begin{omsources}
\item D.~H.~Bailey, J.~M.~Borwein, M.~López de Prado and Q.~J.~Zhu, ``The probability of backtest overfitting'', \emph{Journal of Computational Finance}, 2016.
\item C.~R.~Harvey and Y.~Liu, ``Backtesting'', \emph{Journal of Portfolio Management} 42(1), 2015.
\item M.~López de Prado, \emph{Advances in Financial Machine Learning}, Wiley, 2018; skfolio, \texttt{CombinatorialPurgedCV} (documentation of purging and embargo).
\item R.~Arnott, C.~R.~Harvey and H.~Markowitz, ``A backtesting protocol in the era of machine learning'', \emph{Journal of Financial Data Science} 1(1), 2019.
\end{omsources}

\section{Exercises}

\begin{exercise}[$\star$]\label{exo:rs:overfitting:1}
In a CSCV split the \omterm{def:rs:overfitting:overfitting}{in-sample} winner ranks 380th of 1\,000 out of sample (1 the worst). What is its logit, and does the split count toward the PBO?
\end{exercise}

\begin{exercise}[$\star$]\label{exo:rs:overfitting:2}
How many \omterm{def:rs:vectorised-backtests:backtest}{backtest} paths does CPCV give with 10 folds and 2 test folds? With 6 and 3?
\end{exercise}

\begin{exercise}[$\star$]\label{exo:rs:overfitting:3}
Using the proposition's approximation, how many years does a search over 100 trials need for a target Sharpe ratio of 1?
\end{exercise}

\begin{exercise}[$\star\star$]\label{exo:rs:overfitting:4}
Labels are 10-day forward returns observed daily; the test fold covers days 200 to 299. Which training days does \omterm{def:rs:overfitting:purging}{purging} remove, and which does an \omterm{def:rs:overfitting:purging}{embargo} of
5 days add?
\end{exercise}

\begin{exercise}[$\star\star$]\label{exo:rs:overfitting:5}
Why can the PBO exceed one half? Explain with the thousand correlated trend systems.
\end{exercise}

\begin{exercise}[$\star\star$]\label{exo:rs:overfitting:6}
Why did the planted trend not make the \omterm{def:rs:overfitting:overfitting}{in-sample} winner a good choice, although every system earned the trend?
\end{exercise}

\begin{exercise}[$\star\star\star$]\label{exo:rs:overfitting:7}
\emph{Coding.} Run CSCV on the first 100 of the thousand systems only, and then on 100 independent noise strategies. Compare the PBOs and explain.
\end{exercise}

\begin{exercise}[$\star\star\star$]\label{exo:rs:overfitting:8}
\emph{Find the flaw.} ``We tuned the model's hyperparameters by 10-fold cross-validation with shuffling, which is the standard, so the \omterm{def:rs:overfitting:overfitting}{out-of-sample} estimate is
unbiased.''
\end{exercise}

\section{Problem: One Thousand Trend Systems}

\begin{problem}[{Weekend problem --- a search, measured}]\label{pb:rs:overfitting:1}
The chapter's thousand trend systems on 25 years of synthetic returns, on noise and with a planted trend.

\medskip\noindent\textbf{Part I --- The search.}
\begin{enumerate}
\item What are the best \omterm{def:rs:overfitting:overfitting}{in-sample} Sharpe ratio and its \omterm{def:rs:overfitting:overfitting}{out-of-sample} value on noise? With the trend?
\item What are the median \omterm{def:rs:overfitting:overfitting}{in-sample} and the average \omterm{def:rs:overfitting:overfitting}{out-of-sample} Sharpe ratios in each world?
\item What maximum would a thousand independent null strategies give over 17.9 years, and why is the observed winner lower?
\item What does the trend do to every system, and what does the search add?
\end{enumerate}

\medskip\noindent\textbf{Part II --- The PBO.}
\begin{enumerate}[resume]
\item How does CSCV compute a split's Sharpe ratios quickly?
\item What are the PBOs and the degradation slopes in each world?
\item What does a PBO above one half mean?
\item What would the PBO be for one strategy with a real edge among noise?
\end{enumerate}

\medskip\noindent\textbf{Part III --- Length and leakage.}
\begin{enumerate}[resume]
\item Give the \omterm{def:rs:overfitting:minbtl}{minimum backtest lengths} for a thousand trials at targets of 0.5, 1 and 1.8.
\item Why does halving the target quadruple the length?
\item What $R^2$ do the four cross-validation schemes give for the forecast of noise?
\item Why do shuffled folds leak, and what do \omterm{def:rs:overfitting:purging}{purging} and \omterm{def:rs:overfitting:purging}{embargo} each remove?
\end{enumerate}

\medskip\noindent\textbf{Part IV --- The verdict.}
\begin{enumerate}[resume]
\item What Sharpe ratio does annual walk-forward re-selection give in each world?
\item State the \emph{named result}: the PBO of the search and the deflated \omterm{def:rs:overfitting:overfitting}{out-of-sample} expectation of the chosen system.
\item Which would you trade: the \omterm{def:rs:overfitting:overfitting}{in-sample} winner, the average system, or the walk-forward procedure?
\item What should the \omterm{def:rs:the-research-process:log}{research log} contain for this search?
\item How would a mechanism change your confidence?
\item What does Harvey and Liu's haircut add to the PBO?
\item Where does CPCV improve on walk-forward?
\item In one sentence: what is overfitting in a \omterm{def:rs:vectorised-backtests:backtest}{backtest} search?
\end{enumerate}
\end{problem}

\section{Interview questions}

\begin{interviewq}[$\star$ \iqroles{researcher}]\label{iq:rs:overfitting:1}
You tested 500 strategies and the best has a Sharpe ratio of 2 in the \omterm{def:rs:vectorised-backtests:backtest}{backtest}. What do you do next?
\end{interviewq}

\begin{interviewq}[$\star\star$ \iqroles{researcher, mle}]\label{iq:rs:overfitting:2}
Why is standard k-fold cross-validation wrong for financial time series, and what do you use instead?
\end{interviewq}

\begin{interviewq}[$\star\star$ \iqroles{researcher}]\label{iq:rs:overfitting:3}
Explain the \omterm{def:rs:overfitting:pbo}{probability of backtest overfitting} and how CSCV estimates it.
\end{interviewq}

\begin{interviewq}[$\star\star$ \iqroles{researcher}]\label{iq:rs:overfitting:4}
How long must a \omterm{def:rs:vectorised-backtests:backtest}{backtest} be to trust a Sharpe ratio of 1 found among 1\,000 trials?
\end{interviewq}

\begin{interviewq}[$\star\star$ \iqroles{researcher, risk}]\label{iq:rs:overfitting:5}
What organisational practices reduce overfitting in a research team?
\end{interviewq}

\begin{interviewq}[$\star\star\star$ \iqroles{researcher}]\label{iq:rs:overfitting:6}
Derive the scaling of the \omterm{def:rs:overfitting:minbtl}{minimum backtest length} with the number of trials and the target Sharpe ratio.
\end{interviewq}
